\section{Simulation Bridge Interaction Patterns}
\label{sec:dt_sim_bridge}

The integration of DTs with simulation environments enables a range of interaction patterns designed to support different phases of the system lifecycle and corresponding operational requirements. These interaction patterns are characterized by bidirectional data exchange and control flow, allowing DTs to incorporate predictive capabilities while simulations remain continuously informed by real-time contextual updates.
Building on the DT-SB architecture, we identify three primary interaction paradigms: \emph{Digital Twin \& Batch Simulation}, \emph{Digital Twin \& Streaming Simulation}, and \emph{Physical Twin Simulation}, as illustrated in Figure~\ref{fig:dt_simulation_patterns}. Each paradigm serves distinct objectives, encompassing design validation, what-if scenario analysis, live system adaptation, and virtual commissioning.
Together, they provide a flexible framework for coupling DTs with simulation processes across the development and operational spectrum.


\paragraph{Batch Simulation}
\label{ssec:dt-batch-simulation}

Leverages one-shot simulation runs to analyze scenarios or validate design decisions related to the DT and its physical counterpart. Acting as an intermediary, the DT-SB receives inputs from the DT, such as the current system state, historical data, or scenario configuration parameters and subsequently triggers one or more simulation runs accordingly. These batch simulations execute to completion producing results that the DT-SB forwards back to the DT.
%For instance, a DT of a manufacturing robot may employ batch simulations to evaluate the effects of various arm configurations on cycle time and energy consumption. Similarly, a DT representing a power grid might conduct ``what-if'' analyses to assess the impact of sudden demand surges or component failures on grid stability.
The resulting simulation outcomes enable the DT to refine its models, develop predictive maintenance schedules, and recommend adjustments. This pattern is particularly well suited for use cases requiring comprehensive, offline evaluations, such as design validation and strategic scenario planning.


\paragraph{Streaming Simulation}
\label{ssec:dt-streaming-simulation}

Facilitates dynamic interactions between the DT and the simulation environment through continuous data exchange. In this mode, the DT streams high-frequency operational data to the DT-SB, which relays it to a connected simulation that dynamically adapts its model based on the incoming input. The simulation continuously produces output data that is streamed back to the DT through the DT-SB.
This interaction enables live behavioral analysis and adaptive control.
%For example, a DT of a wind turbine can stream sensor data to a simulation predicting fatigue life under current wind loads; simulation results subsequently inform real-time pitch adjustments to optimize turbine health. Similarly, in smart buildings, a DT can supply HVAC sensor data to an energy optimization simulation, enabling real-time control commands to minimize energy consumption based on predicted occupancy and weather conditions.
Furthermore, deviations between live DT data and simulation outputs can serve as early indicators of anomalies or faults, thereby enabling proactive intervention.
This pattern is particularly well suited for applications demanding immediate feedback and continuous system optimization.


\paragraph{Physical Twin Simulation}
\label{ssec:physical-twin-simulation}

Involves the use of a simulation environment that functions as a digital replica of the PT.
%a physical entity or environment often referred to as a \emph{Physical Twin} (PT).
This interaction mode is particularly valuable for virtual commissioning, operator training, and control logic testing, where system behaviors can be evaluated in a safe, simulated setting prior to deployment in the real world.
In this pattern, the DT is unaware of the fact that it is communicating with a simulator and exchange data and control commands that are actuated in the simulated environment.
DT developers can configure simulation parameters through the DT-SB in order to test the behavior of the DT under varying conditions.
% In this pattern, the DT interacts with the simulator through the DT-SB, configuring simulation parameters, initiating test scenarios, and issuing control commands. The simulator, in turn, provides logs, metrics, and state feedback to the DT, effectively representing the outcomes of these interactions.
%
% For instance, prior to deploying a new production line, its DT may control a simulated model of the line to evaluate control algorithms, identify bottlenecks, and optimize throughput—without risking downtime or production losses.
% Similarly, operators can train in DT-driven virtual environments, rehearsing complex procedures or emergency protocols without exposing actual equipment to risk.
This pattern enables thorough testing and validation within a controlled virtual environment, significantly reducing the costs and risks associated with physical prototyping and on-site experimentation.
